\documentclass[10pt,a4paper]{amsart}
\usepackage[margin=19mm]{geometry}
\usepackage{amsmath,amssymb,mathtools}
\usepackage[scr=rsfs,cal=euler]{mathalfa}
\usepackage{xfrac}
\usepackage[unicode,hidelinks,bookmarks=false]{hyperref}
\newcommand{\ie}{i.e.\ }
\newcommand{\cf}{cf.\ }
\newcommand{\resp}{respectively\ }
\newcommand{\Subsection}{\subsection}
\providecommand{\nobreakdash}{\nobreak\hbox{-}}
\setlength{\emergencystretch}{2em}
\multlinegap0pt
\usepackage[shortlabels]{enumitem}
\usepackage{xcolor}
\usepackage{amssymb}
\usepackage{csquotes}
\usepackage{algorithm}\usepackage{algpseudocode}
\usepackage{tikz}
\usepackage{mathtools}
\newtheorem{assumption}{Assumption}
\newtheorem{theorem}{Theorem}
\title{Policy Optimization in Hybrid Discrete-Continuous Action Spaces via Mixed Gradients}
\author{Matias Alvo, Daniel Russo, Yash Kanoria}
\date{Source: arXiv:2605.14297v1 (CC BY 4.0)}
\begin{document}
\maketitle

\noindent\emph{Source and licence.} Policy Optimization in Hybrid Discrete-Continuous Action Spaces via Mixed Gradients. Version arXiv:2605.14297v1 (CC BY 4.0), distributed by arXiv under the Creative Commons Attribution 4.0 International licence, http://creativecommons.org/licenses/by/4.0/, as recorded in the arXiv licence metadata for this exact version and shown on the abstract page https://arxiv.org/abs/2605.14297v1. Full licence notice: \texttt{licenses/arxiv-2605.14297-CC-BY-4.0.txt}.\par\smallskip

\noindent\textit{Editorial scope.} Counted unit: the article's theorem thm:cross-term-vanishes-for-near-opt (source0.tex lines 626--640). The article's complete original proof of it is delivered with it (source0.tex lines 1244--1490, one \texttt{proof} environment of 247 lines, which the article places away from the statement). No mathematics is written, repaired or re-derived: the statement and the proof are the article's own lines, in the article's own notation, and no retained line is substituted or re-flowed.  Retained uncounted as the source-local context the proof needs: lines 613--623.  Nothing was omitted from any retained span, so the package carries no omission record.  No retained line contains a citation, so the wrapper carries no bibliography and no bibliography entry.  No retained line contains a figure environment, an image-inclusion command or drawing code, so no image is delivered and the package declares no asset dependency.  Editorial formatting only, in addition: an AMS \texttt{amsart} preamble that restates the article's own declarations and macros used by the retained lines, namely the environments \texttt{assumption}, \texttt{theorem} on separate counters, together with the standard packages those lines need.  The retained environments print in this document as Assumption 1, Theorem 1 in order of appearance, whereas the article prints the same items with the numbering of its own full document.  The complete native source of the article as distributed in the arXiv e-print for this exact version is bundled unchanged as \texttt{sources/200572/source0.tex}.
\begin{assumption}[Regularity conditions]
\label{ass:regularity-cross-term}
There exist constants $C, L_\pi, L_s < \infty$ such that, for all $\kappa \in \mathcal{K}$:
\begin{enumerate}[leftmargin=*, noitemsep]
    \item \textbf{Statewise best response.} There exists best-response $\pi^{\mathcal{X}}(\kappa)$ such that for all $s \in \mathcal{S}$, $V_{\pi^*_\kappa}(s) = \min_{\pi^{\mathcal{X}}} V_{(\pi^{\mathcal{X}}, \pi^{\mathcal{B}}_\kappa)}(s)$.
    \item \textbf{Uniform near-best response.} If $\pi^{\mathcal{X}}_\phi$ is %an 
    any $\epsilon$-best response to $\kappa$, then for any initial state distribution $\rho'$, $\mathbb{E}_{s_0 \sim \rho'}[V_{\pi_{\theta}}(s_0) - V_{\pi^*_\kappa}(s_0)] \leq C \epsilon$.
    \item \textbf{Policy smoothness and state sensitivity.} 
    $\|\nabla_\kappa s_t(\kappa)\| \leq L_s$ for all $t\geq 1$, $\xi_{0:t-1}$ and $x_{0:t-1} \in \mathcal X^{t}$, and $\sup_{\phi, s, x} \|\nabla_s \log \pi^{\mathcal{X}}_\phi(x \mid s)\| \leq L_\pi$.
\end{enumerate}
\end{assumption}

\begin{theorem}[Cross term vanishes for near-best responses]
\label{thm:cross-term-vanishes-for-near-opt}
Under Assumption~\ref{ass:regularity-cross-term}, for any $\kappa$ and any $\epsilon$-best response $\pi^{\mathcal{X}}_\phi$,
\[
\left\|
\mathbb{E}_{\theta}\!\left[
\sum_{t\geq 0}
\gamma^{t}\,
\nabla_{\kappa}\log \pi^{\mathcal{X}}_{\phi}\!\left(x_t \mid s_t(\kappa)\right)\,
Q_{\pi_{\theta}}(s_t(\kappa),a_t)
\right]
\right\|
\;\leq \epsilon  \times  L_{\pi} L_s  \left(1 + \frac{C \gamma}{1 - \gamma} \right).
\]
\end{theorem}

\begin{proof}[Proof of Theorem~\ref{thm:cross-term-vanishes-for-near-opt}]

For a fixed continuous policy parameter $\kappa$, define the post-discrete
action-value function
\[
\tilde Q_{\pi}(s,x)
:=
Q_{\pi}\bigl(s,(x,\pi^{\mathcal B}_{\kappa}(s, x))\bigr),
\]
that is, the value of selecting discrete action $x$ and then applying the
continuous action prescribed by $\pi^{\mathcal B}_{\kappa}$ conditional on $x$. Similarly,
define the post-discrete advantage of the best-response policy $\pi^*_\kappa$ as
\[
\tilde A_{\pi^*_\kappa}(s,x)
:=
\tilde Q_{\pi^*_\kappa}(s,x)-V_{\pi^*_\kappa}(s).
\]
Since $\pi^*_\kappa$ is a statewise best response by
Assumption~\ref{ass:regularity-cross-term}(1), we have
\[
\tilde A_{\pi^*_\kappa}(s,x) \geq 0
\qquad \text{for all } (s,x).
\]

By the performance difference lemma for costs,
\[
\mathbb{E}_{s_0\sim \rho}
\!\left[
V_{\pi_\theta}(s_0)-V_{\pi^*_\kappa}(s_0)
\right]
=
\mathbb{E}_{\theta}\!\left[
\sum_{t\geq 0}\gamma^t
\tilde A_{\pi^*_\kappa}(s_t(\kappa),x_t)
\right].
\]
Since $\pi^{\mathcal X}_{\phi}$ is an $\epsilon$-best response to
$\kappa$, the left-hand side is at most $\epsilon$, and therefore
\begin{equation}
\label{eq:advantage-bound}
\mathbb{E}_{\theta}\!\left[
\sum_{t\geq 0}\gamma^t
\tilde A_{\pi^*_\kappa}(s_t(\kappa),x_t)
\right]
\leq \epsilon.
\end{equation}

Let
\[
\mathcal C
:=
\mathbb{E}_{\theta}\!\left[
\sum_{t\geq 0}
\gamma^t
\nabla_{\kappa}\log \pi^{\mathcal X}_{\phi}
\!\left(x_t\mid s_t(\kappa)\right)
Q_{\pi_\theta}(s_t(\kappa),a_t)
\right]
\]
denote the cross term. Since
\[
Q_{\pi_\theta}(s_t(\kappa),a_t)
=
\tilde Q_{\pi_\theta}(s_t(\kappa),x_t),
\]
we can write
\begin{align*}
\mathcal C
&=
\mathbb{E}_{\theta}\!\left[
\sum_{t\geq 0}\gamma^t
\nabla_{\kappa}\log \pi^{\mathcal X}_{\phi}
\!\left(x_t\mid s_t(\kappa)\right)
\tilde Q_{\pi_\theta}(s_t(\kappa),x_t)
\right] \\
&=
\mathbb{E}_{\theta}\!\left[
\sum_{t\geq 0}\gamma^t
\nabla_{\kappa}\log \pi^{\mathcal X}_{\phi}
\!\left(x_t\mid s_t(\kappa)\right)
\tilde Q_{\pi^*_\kappa}(s_t(\kappa),x_t)
\right] \\
&\quad+
\mathbb{E}_{\theta}\!\left[
\sum_{t\geq 0}\gamma^t
\nabla_{\kappa}\log \pi^{\mathcal X}_{\phi}
\!\left(x_t\mid s_t(\kappa)\right)
\left(
\tilde Q_{\pi_\theta}(s_t(\kappa),x_t)
-
\tilde Q_{\pi^*_\kappa}(s_t(\kappa),x_t)
\right)
\right].
\end{align*}

For fixed realized history up to time $t$ excluding $x_t$, the state
$s_t(\kappa)$ is fixed as a differentiable function of $\kappa$. Hence
\[
\sum_{x\in\mathcal X}
\pi^{\mathcal X}_\phi(x\mid s_t(\kappa))
\nabla_\kappa \log \pi^{\mathcal X}_\phi(x\mid s_t(\kappa))
=
\nabla_\kappa
\sum_{x\in\mathcal X}
\pi^{\mathcal X}_\phi(x\mid s_t(\kappa))
=0.
\]
Thus, any term depending only on $s_t(\kappa)$ may be subtracted from
$\tilde Q_{\pi^*_\kappa}(s_t(\kappa),x_t)$ inside the conditional expectation over $x_t$.
Therefore,
\begin{align*}
\mathcal C
&=
T_1+T_2,
\end{align*}
where
\[
T_1
:=
\mathbb{E}_{\theta}\!\left[
\sum_{t\geq 0}\gamma^t
\nabla_{\kappa}\log \pi^{\mathcal X}_{\phi}
\!\left(x_t\mid s_t(\kappa)\right)
\tilde A_{\pi^*_\kappa}(s_t(\kappa),x_t)
\right],
\]
and
\[
T_2
:=
\mathbb{E}_{\theta}\!\left[
\sum_{t\geq 0}\gamma^t
\nabla_{\kappa}\log \pi^{\mathcal X}_{\phi}
\!\left(x_t\mid s_t(\kappa)\right)
\left(
\tilde Q_{\pi_\theta}(s_t(\kappa),x_t)
-
\tilde Q_{\pi^*_\kappa}(s_t(\kappa),x_t)
\right)
\right].
\]

By Assumption~\ref{ass:regularity-cross-term}(3),
\[
\left\|
\nabla_{\kappa}\log \pi^{\mathcal X}_{\phi}
(x_t\mid s_t(\kappa))
\right\|
\leq
\left\|
\nabla_s \log \pi^{\mathcal X}_{\phi}
(x_t\mid s_t(\kappa))
\right\|
\left\|
\nabla_\kappa s_t(\kappa)
\right\|
\leq
L_\pi L_s.
\]
Using this bound and \eqref{eq:advantage-bound},
\begin{align*}
\|T_1\|
&\leq
\mathbb{E}_{\theta}\!\left[
\sum_{t\geq 0}\gamma^t
\left\|
\nabla_{\kappa}\log \pi^{\mathcal X}_{\phi}
(x_t\mid s_t(\kappa))
\right\|
\tilde A_{\pi^*_\kappa}(s_t(\kappa),x_t)
\right] \\
&\leq
L_\pi L_s\,
\mathbb{E}_{\theta}\!\left[
\sum_{t\geq 0}\gamma^t
\tilde A_{\pi^*_\kappa}(s_t(\kappa),x_t)
\right] \\
&\leq
L_\pi L_s\,\epsilon.
\end{align*}

It remains to bound $T_2$. For any $(s_t,x_t)$, both
$\tilde Q_{\pi_\theta}(s_t(\kappa),x_t)$ and $\tilde Q_{\pi^*_\kappa}(s_t(\kappa),x_t)$ use the
same immediate cost and the same continuous action
$\pi^{\mathcal B}_{\kappa}(s_t(\kappa),x_t)$ in the first period. They differ
only in the continuation policy. Hence,
\begin{align}
\label{eq:q-difference-bound}
\tilde Q_{\pi_\theta}(s_t(\kappa),x_t)
-
\tilde Q_{\pi^*_\kappa}(s_t(\kappa),x_t)
&=
\gamma\,
\mathbb{E}\!\left[
V_{\pi_\theta}(s_{t+1})
-
V_{\pi^*_\kappa}(s_{t+1})
\mid s_t(\kappa),x_t
\right] \\
&\leq
\gamma C\epsilon,
\nonumber
\end{align}
where the inequality follows from Assumption~\ref{ass:regularity-cross-term}(2)
applied to the conditional distribution of $s_{t+1}$ given $(s_t(\kappa),x_t)$.
Moreover, the left-hand side is nonnegative by the statewise optimality of
$\pi^*_\kappa$.

Using \eqref{eq:q-difference-bound} and Assumption~\ref{ass:regularity-cross-term}(3),
\begin{align*}
\|T_2\|
&\leq
\mathbb{E}_{\theta}\!\left[
\sum_{t\geq 0}\gamma^t
\left\|
\nabla_{\kappa}\log \pi^{\mathcal X}_{\phi}
(x_t\mid s_t(\kappa))
\right\|
\left|
\tilde Q_{\pi_\theta}(s_t(\kappa),x_t)
-
\tilde Q_{\pi^*_\kappa}(s_t(\kappa),x_t)
\right|
\right] \\
&\leq
L_\pi L_s
\mathbb{E}_{\theta}\!\left[
\sum_{t\geq 0}\gamma^t
\gamma C\epsilon
\right] \\
&=
L_\pi L_s\,
\frac{C\gamma}{1-\gamma}\,
\epsilon.
\end{align*}

Combining the bounds for $T_1$ and $T_2$ gives
\[
\|\mathcal C\|
\leq
L_\pi L_s\,\epsilon
\left(
1+\frac{C\gamma}{1-\gamma}
\right),
\]
which proves the result.
\end{proof}

\end{document}
